% ============================================================
% V. ITEM 4 — PRECISIÓN DE MUESTREO EN ESP32-WROOM
% ============================================================

\section{Precisión de muestreo en ESP32-WROOM}
\label{sec:embebidos}

Como se describe en la Sección~\ref{sec:metodologia-embebidos}, esta sección se basa en una \textbf{simulación etiquetada}, generada por \texttt{firmware/simulate\_serial\_output.py} a partir de un modelo de temporización explícito (velocidad de suma $\approx 2\times10^{7}$ sumas/s, \emph{jitter} de FreeRTOS $\pm2$ ms, \emph{jitter} de temporizador de hardware $\pm1$ ms), y \textbf{no} de una captura de puerto serie real, al no disponerse de un ESP32-WROOM físico para flashear los tres sketches de \texttt{firmware/}.

\begin{table}[!t]
\centering
\caption{Periodo de muestreo simulado promedio (8 muestras) por mecanismo}
\label{tab:embebidos}
\small
\begin{tabular}{@{}l r r@{}}
\toprule
Mecanismo & $N$ (suma) & $\overline{dt}$ (ms) \\
\midrule
4(a) bucle bloqueante & 1\,000\,000 & 150.0 \\
4(a) bucle bloqueante & 10\,000\,000 & 599.9 \\
4(a) bucle bloqueante & 100\,000\,000 & 5100.1 \\
\midrule
4(b) doble núcleo & 1\,000\,000 & 100.2 \\
4(b) doble núcleo & 10\,000\,000 & 99.6 \\
4(b) doble núcleo & 100\,000\,000 & 99.8 \\
\midrule
4(c) interrupción & 1\,000\,000 & 100.4 \\
4(c) interrupción & 10\,000\,000 & 150.0 \\
4(c) interrupción & 100\,000\,000 & 712.7 \\
\bottomrule
\end{tabular}
\end{table}

\subsection{Efecto sobre la frecuencia de muestreo}

El mecanismo 4(a) es el único de los tres en el que el periodo de muestreo efectivo escala de forma prácticamente lineal con $N$: pasa de $\sim\!150$ ms ($N=10^6$) a $\sim\!5.1$ s ($N=10^8$), muy lejos del periodo nominal de $100$ ms, porque la suma y el muestreo comparten el mismo hilo bloqueante. El mecanismo 4(b) mantiene el periodo nominal ($\sim\!100$ ms $\pm$ \emph{jitter} de planificación) prácticamente constante para los tres valores de $N$, porque la tarea de muestreo corre en un núcleo físicamente distinto al de la suma. El mecanismo 4(c) mantiene la referencia temporal exacta del temporizador de hardware, pero el promedio de las muestras \emph{procesadas} por \texttt{loop()} se ve inflado para $N$ grande por la muestra ``atrasada'' que acumula el retraso de procesamiento —una distinción importante entre \emph{cuándo ocurre la interrupción} (exacto) y \emph{cuándo \texttt{loop()} logra atenderla} (retrasado si la suma es larga).

\subsection{¿Cuál mecanismo es más preciso?}

Con base en el modelo simulado, el orden de precisión del \emph{periodo de muestreo alcanzado} (no solo de la referencia temporal) es:

\begin{equation}
\text{4(b) doble núcleo} \;>\; \text{4(c) interrupción} \;>\; \text{4(a) bucle bloqueante}
\end{equation}

El mecanismo de \textbf{doble núcleo (4b)} es el más preciso en términos prácticos: al ejecutar el muestreo y la carga de cómputo en núcleos físicamente separados, el periodo de muestreo efectivo permanece cercano al nominal independientemente de $N$. El mecanismo de \textbf{interrupción (4c)} ofrece la referencia temporal más exacta de los tres a nivel de hardware —el temporizador dispara exactamente cada $100$ ms sin importar qué esté haciendo \texttt{loop()}—, pero como el propio \texttt{loop()} debe procesar cada muestra (leer los potenciómetros e imprimir), una suma suficientemente larga ($N=10^8$) provoca que el procesamiento se atrase y varias interrupciones se acumulen antes de ser atendidas, degradando la precisión \emph{observada} aunque no la del disparo del temporizador en sí. El \textbf{bucle bloqueante (4a)} es, sin ambigüedad, el menos preciso: la suma bloqueante retrasa directamente cada muestra, de forma proporcional a $N$, sin ningún mecanismo que lo evite.

En síntesis: si el requisito es una frecuencia de muestreo constante y predecible independientemente de la carga de cómputo, el mecanismo de doble núcleo es la opción más robusta; si el requisito es una referencia temporal exacta con la posibilidad de procesar cada evento algo más tarde (tolerando cierta latencia de procesamiento, no de disparo), la interrupción de hardware es preferible sobre el bucle bloqueante, y ambas alternativas evitan por completo la degradación lineal que sufre el mecanismo de un solo hilo.
